\documentclass[a4paper]{article}

\usepackage{graphicx}
\usepackage{amsmath,amssymb}
\usepackage{minted}
\usepackage{multirow}
\usepackage{float}
\usepackage{algorithm}
\usepackage{algpseudocode}
\usepackage{todonotes}
\usepackage{forest}
\usepackage{xstring}
\usepackage{xcolor}
\usepackage[most]{tcolorbox}
\usepackage{xparse}
\usepackage{natbib}
\usepackage{adjustbox}

\usetikzlibrary{graphs,graphdrawing}
\usegdlibrary{layered}

% for external graphviz
\input{/home/donkarlo/Dropbox/repo/nd_language_ai_project/src/nd_language_ai/formal/computer/programming/tex/latex/code_clips/external_graphviz.tex}

% for tags
\input{/home/donkarlo/Dropbox/repo/nd_language_ai_project/src/nd_language_ai/formal/computer/programming/tex/latex/part/control_sequence/kind/semtag/semtag.tex}

% for links
\input{/home/donkarlo/Dropbox/repo/nd_language_ai_project/src/nd_language_ai/formal/computer/programming/tex/latex/code_clips/seehere.tex}

% for nested lists and sections
\input{/home/donkarlo/Dropbox/repo/nd_language_ai_project/src/nd_language_ai/formal/computer/programming/tex/latex/code_clips/deep_structure.tex}

\title{Auto regression}
\date{}
\author{Mohammad Rahmani}

\begin{document}
   \maketitle

   \begin{table}[H]

      \centering

      \small

      \setlength{\tabcolsep}{4pt}

      \begin{tabular}{p{0.25\textwidth} p{0.18\textwidth} p{0.47\textwidth}}

         \hline

         \textbf{Concept} & \textbf{Variable name} & \textbf{Description} \\

         \hline

         Spike-train feature vector
         & $\mathbf{q}_n$
         & Vector containing the observation-neuron firing rates
         $r_{Z_i^{(d)}}([t_n,t_{n+1}])$ at discrete time index $n$. \\

         Discrete time index
         & $n$
         & Index of a firing-rate observation interval in the spike-train sequence. \\

         GRU input-window length
         & $W$
         & Number of consecutive spike-train feature vectors contained in each overlapping input window. \\

         GRU encoder
         & $\mathcal{E}_{\phi}$
         & GRU-based encoder that maps an input window of spike-train feature vectors to a latent sequence. \\

         GRU decoder
         & $\mathcal{D}_{\psi}$
         & Decoder that reconstructs the original spike-train feature window from its latent representation. \\

         Encoder parameters
         & $\phi$
         & Trainable parameters of the GRU encoder. \\

         Decoder parameters
         & $\psi$
         & Trainable parameters of the decoder. \\

         Latent dimension
         & $K$
         & Number of features contained in each latent representation produced by the GRU encoder. \\

         Latent spike-train representation
         & $\mathbf{h}_n$
         & $K$-dimensional latent representation associated with discrete time index $n$, obtained by averaging overlapping GRU encodings. \\

         Latent-feature index
         & $k$
         & Index of a component of the latent representation, where $k \in \{1,\ldots,K\}$. \\

         Sequence length
         & $N_{\mathrm{GRU}}$
         & Number of latent time steps used for autocorrelation-based period discovery. \\

         Autocorrelation lag
         & $\tau$
         & Temporal displacement, measured in discrete time steps, between two latent spike-train representations. \\

         Latent autocorrelation
         & $\rho_{\mathrm{GRU}}(\tau)$
         & Multivariate autocorrelation of the GRU latent representation at lag $\tau$, computed efficiently using the Fast Fourier Transform. \\

         Minimum admissible period
         & $\tau_{\min}$
         & Smallest autocorrelation lag considered as a possible period. \\

         Estimated period
         & $\hat{\tau}_{\mathrm{per}}$
         & First detected local maximum of $\rho_{\mathrm{GRU}}(\tau)$ for $\tau \geq \tau_{\min}$. \\

         Fast Fourier Transform
         & $\operatorname{FFT}(\cdot)$
         & Transformation used to compute the latent autocorrelation efficiently in the frequency domain. \\

         Inverse Fast Fourier Transform
         & $\operatorname{IFFT}(\cdot)$
         & Inverse transformation used to convert the frequency-domain autocorrelation product back to the lag domain. \\

         \hline

      \end{tabular}

      \caption{Variables used in the GRU-based latent autocorrelation period-discovery method.}

      \label{tab:gru-fft-period-discovery-variables}

   \end{table}

   \bibliographystyle{plainnat}
   \bibliography{/home/donkarlo/Dropbox/repo/nd_shared_research/bibliography/bibliography.bib}
\end{document}
